\section{Appendix}

\subsection{Related Work}

This section reviews the work most closely related to the problem studied in this paper. Since our goal is not to propose a new intrusion detector or a new attack simulation platform, we focus on three lines of work that are directly connected to the limitations discussed in the introduction: attack knowledge bases and product evaluations, defender-side knowledge bases, and quantitative risk or attack-path analysis.

\subsubsection{Attack Knowledge Bases and Step-Level Product Evaluation}

MITRE ATT\&CK has become the most widely used knowledge base for describing adversary behavior. It organizes tactics, techniques, and procedures observed in real-world intrusions and provides a common language for threat intelligence, red-team exercises, and security product evaluation~\cite{mitreAttack}. In practical terms, ATT\&CK answers the question: \emph{what did the attacker do?} This is an important starting point for any APT-related evaluation, because without a shared vocabulary of attacker behavior, different reports and tests are difficult to compare.

However, the question addressed in this paper is different. Knowing that an attacker used a certain ATT\&CK technique does not by itself explain why the victim environment failed to stop it. The same technique may succeed because of weak password policies, missing multi-factor authentication, exposed remote services, weak execution control, insufficient data protection, or poor monitoring. These causes lead to different remediation actions, even when the attacker behavior is described by the same ATT\&CK technique. Therefore, ATT\&CK provides the behavioral description of the attack, but it does not provide a defender-side failure attribution.

MITRE ATT\&CK Evaluations move one step closer to defense assessment by reporting how evaluated products detect or protect against emulated attack steps~\cite{mitreAttackEvaluations}. This line of work answers a second question: \emph{did the product observe, detect, or block this step?} Such step-level results are useful, especially when comparing product visibility and protection behavior under a shared emulation plan. Still, they mainly produce a visibility and protection matrix. A missed or delayed detection is reported as an outcome, but the result does not explain which victim-side control point first failed. It also does not directly distinguish whether the weakness belongs to identity control, network isolation, endpoint execution protection, data protection, or observability.

This distinction is important for our work. A step-level matrix can tell defenders where a product responded, but it does not tell them how to attribute failures to repairable control points. Our framework uses ATT\&CK-style attack descriptions as input, but it adds a weakness attribution layer on the defender side.

\subsubsection{Defensive Knowledge Bases and Control Capabilities}

D3FEND complements ATT\&CK by organizing cybersecurity countermeasure techniques from a defender-oriented perspective~\cite{mitreD3FEND,kaloroumakis2021d3fend}. If ATT\&CK asks what the attacker did, D3FEND asks a different question: \emph{if this defensive technique is deployed, what adversary behavior or digital artifact can it counter?} This makes D3FEND useful for reasoning about defensive capabilities and for connecting adversary behaviors to possible countermeasures.

Nevertheless, D3FEND is still not a scoring model for multi-stage APT defense. It describes defensive techniques and their relationships to artifacts and adversary behaviors, but it does not decide which control point failed first in a particular attack action. It also does not define how the result of one action should be combined with other actions in the same stage, or how stage-level results should contribute to an overall playbook score.

In our setting, this difference matters. A defense knowledge base such as D3FEND can indicate what kinds of techniques may counter an attack behavior, but a scoring framework must answer a more operational question: \emph{which victim-side control point failed in this action, and how should this failure affect the score of the defense system?} Our control-point attribution schema and defense-technology crosswalk are designed for that purpose. They do not replace D3FEND; rather, they use a complementary viewpoint centered on control-point failure and repair priority.
\subsubsection{Vulnerability Scoring and Attack-Path Analysis}

A third related line of work focuses on quantifying vulnerability severity and attack-path risk. CVSS provides a standardized way to measure the severity of software vulnerabilities~\cite{mell2007common,first2019cvss31}. It answers the question: \emph{how severe is this vulnerability?} This is useful for vulnerability management, but it is not designed to evaluate how a deployed defense system performs against a multi-stage APT playbook. A vulnerability may have a high severity score, but that score does not tell us whether the defense system blocks the initial access, detects the execution step, prevents lateral movement, or only observes the attack near the exfiltration stage.

Attack-graph research extends the analysis from individual vulnerabilities to multi-step attack paths. MulVAL and related work model how vulnerabilities and configurations can be combined into reachable attack paths~\cite{ou2005mulval,ou2006scalable}. Other studies further introduce quantitative or probabilistic security metrics over attack graphs~\cite{wang2007measuring,wang2008probabilistic,frigault2008dynamic}. These methods answer questions such as: \emph{which attack paths are reachable?} and \emph{how risky is a path or network configuration?} More recent cyber risk quantification work also estimates potential loss under different threat and control assumptions~\cite{erola2022system}.

These approaches are valuable, but their unit of analysis is usually a vulnerability, a configuration dependency, an attack path, or a loss distribution. Our unit of analysis is an APT action mapped to a victim-side defensive weakness and then aggregated through stage and playbook logic. In other words, attack graphs mainly quantify attacker reachability or path risk, whereas our framework quantifies defense performance over a given set of APT playbooks.

\subsubsection{Position of This Work}

The above work provides important foundations, but each line answers a different question. ATT\&CK answers what the attacker did. ATT\&CK Evaluations answer whether a product observed, detected, or blocked a step. D3FEND answers what a defensive technique can counter. CVSS and attack graphs answer how severe a vulnerability is or how risky an attack path may be.

This paper addresses a different question: \emph{which victim-side control point first failed, and how should that failure contribute to a stage-weighted defense score?} To answer it, we build a control-point attribution schema, map APT actions to first-failed victim-side control points, connect these control points to defense product categories, and aggregate action-level results into stage-level and playbook-level scores. This design directly addresses black-box scoring by producing decomposable scores, addresses undifferentiated stage value through phase weights, and addresses missing attribution by mapping each action to a repairable victim-side control-point failure.






